\chapter{Justificación}\label{chapter:PA}

La fotopletismografía remota (rPPG) tiene el potencial de revolucionar el monitoreo del estrés clínico mediante el análisis no invasivo de la Variabilidad de la 
Frecuencia Cardíaca (HRV). Sin embargo, es de vital importancia abordar el impacto de la iluminación deficiente porque el éxito actual de estos sistemas está 
fuertemente sesgado hacia condiciones ideales, impidiendo su uso masivo. Al diseñar e implementar métodos de rPPG, la gran mayoría de los investigadores asumen 
condiciones operativas ideales, pasando por alto las enormes complejidades de desplegar estos sistemas en entornos del mundo real \cite{casado2024design}. 
Z. Yang y otros confirman este sesgo metodológico al señalar que los estudios existentes exploran casi exclusivamente condiciones de laboratorio con buena 
iluminación, ignorando que en los espacios interiores reales la distribución de la luz es irregular y sufre variaciones frecuentes \cite{yang2022assessment_dl}. 
Esta dependencia de la luz perfecta evidencia un vacío crítico y demuestra la necesidad urgente de evaluar esta tecnología en escenarios realistas y cotidianos 
donde la iluminación fluctúa drásticamente.\\

El impacto de la oscuridad en la precisión del diagnóstico es un problema que no se ha logrado resolver mediante preprocesamiento visual clásico. Al evaluar 
arquitecturas profundas (como DeepPhys y PhysFormer) frente a videos con iluminación reducida, S. Chen, K. Wong y otros demostraron en sus métricas que aplicar filtros 
tradicionales de mejora de imagen (como la Ecualización de Histograma y LIME) en realidad empeora el Error Absoluto Medio (MAE) y destruye la señal rPPG \cite{chen2024image_en}. Peor 
aún, la fiabilidad clínica se rompe precisamente cuando el paciente experimenta estrés en la oscuridad. En sus ensayos experimentales, B. Acharya y W Saakyan comprobaron 
a nivel práctico que cuando se evalúan los modelos ante la combinación simultánea de un ritmo cardíaco elevado (estrés) y baja iluminación (oscuridad), los errores 
de estimación de las redes se disparan drásticamente, demostrando que ningún método actual logra un rendimiento consistente bajo estas dos presiones simultáneas \cite{acharya2025reliability}.

Frente a esta desventaja óptica, es ideal encontrar una alternativa de robustez entre las arquitecturas de extracción local, como podrían ser las Redes 
Neuronales Convolucionales (CNN), y las de atención global, como los Vision Transformers (ViT), ya que los experimentos recientes revelan fallas estructurales en 
ambas topologías al enfrentar ruido \cite{chang2025physvim}. Por un lado, al desarrollar sus nuevos modelos de procesamiento temporal como PhysVim, donde Z. Chang y otros, identificaron a nivel práctico que la 
principal limitación de las CNN por si solas es su campo receptivo local, una debilidad arquitectónica que les impide modelar dependencias espacio-temporales largas y 
provoca fallos sistémicos cuando la imagen presenta ruido. Aunque los ViT prometen una atención sobre todo el rostro, los experimentos de 
N. Nguyen, L. Nguyen y otros, demostraron su fragilidad en entornos complejos: al someter arquitecturas ViT (como EfficientPhys y ContrastPhys) a oclusión, pérdida de calidad, 
compresión de bits y ruido espacial sintético, sus tablas de resultados evidenciaron un colapso drástico en la precisión de los Transformers frente a la 
degradación visual \cite{nguyen2024artifact_resilience}.\\

Aportar a la resolución de este problema otorgará beneficios directos a la ingeniería biomédica al responder de qué manera afecta la iluminación deficiente a la 
precisión del diagnóstico de estrés. Los hallazgos establecerán una frontera medible que determine en qué punto la oscuridad contamina la señal rPPG y 
altera el HRV. Esto es una necesidad ya que los experimentos en laboratorio de Z. Yang, H. Wang y F. Lu demuestran que, aunque las redes profundas funcionan bien con 
buena luz, tienden a colapsar debido a la falta datasets con objetivos específicos pero diferente iluminación \cite{yang2022assessment_dl}. Operar con base a está degradación desde el principio
evitará el efecto cascada, impidiendo que una mala señal óptica genere un falso diagnóstico psicológico.\\

% DESPUÉS
Este proyecto de investigación busca contribuir a la exploración de sensores no invasivos, accesibles y económicos mediante el uso de tecnologías de fotopletismografía
remota (rPPG). Al identificar y evaluar qué arquitecturas de aprendizaje profundo ofrecen mayor resistencia ante condiciones de iluminación baja, el 
estudio pretende facilitar la detección temprana del estrés, permitiendo intervenciones preventivas antes de que se manifiesten trastornos crónicos de 
salud mental.